\chapter{Capitolo 5 - Isomorfismi, equivalenze, limiti e colimiti con vincolo ordinativo}

\section{1. Problema}

Il Capitolo 4 ha definito l'estensione ordinativa delle unità di base: oggetti, morfismi, composizione e identità. Questo è sufficiente per controllare la dinamica locale di una categoria ordinativa, ma non basta per affrontare i costrutti classici che, nella pratica matematica, determinano la potenza reale del formalismo:

\begin{enumerate}
\def\labelenumi{\arabic{enumi}.}
\tightlist
\item
  isomorfismi;
\item
  equivalenze;
\item
  limiti;
\item
  colimiti.
\end{enumerate}

Nel quadro classico questi costrutti sono centrali perché esprimono nozioni di ``stessa struttura'', ``stessa teoria fino a traduzione'', ``migliore oggetto universale di compatibilità'' e ``migliore oggetto universale di aggregazione''. In OCT non vengono rimossi; vengono sottoposti a un vincolo aggiuntivo: validità ordinativa nel contesto \texttt{Omega}.

Il problema del capitolo è quindi:

\textbf{come preservare l'apparato universale classico senza trattare come equivalenti o universalmente validi costrutti che risultano sterili o degenerativi in termini ordinativi?}

La difficoltà non è banale:

\begin{enumerate}
\def\labelenumi{\arabic{enumi}.}
\tightlist
\item
  se il filtro ordinativo è troppo rigido, può distruggere la trasferibilità dei risultati;
\item
  se è troppo permissivo, non evita gli errori che OCT vuole correggere.
\end{enumerate}

Per questo il capitolo adotta una strategia di ``selettività controllata'':

\begin{enumerate}
\def\labelenumi{\arabic{enumi}.}
\tightlist
\item
  la struttura classica resta necessaria;
\item
  la validità ordinativa decide l'operatività reale del costrutto nel dominio osservato.
\end{enumerate}

\subsection{1.1 Nodo critico: ``stesso'' non significa sempre ``equivalente operativamente''}

Nel classico:

\begin{enumerate}
\def\labelenumi{\arabic{enumi}.}
\tightlist
\item
  isomorfismo implica sostituibilità strutturale;
\item
  equivalenza categoriale implica stessa teoria ``a meno di traduzione''.
\end{enumerate}

In OCT:

\begin{enumerate}
\def\labelenumi{\arabic{enumi}.}
\tightlist
\item
  queste nozioni restano valide sul piano sintattico;
\item
  non implicano automaticamente equivalenza ordinativa su funzione, persistenza e robustezza.
\end{enumerate}

La differenza è la stessa già incontrata nei capitoli precedenti, ma qui viene formalizzata sui costrutti più forti del linguaggio categoriale.

\subsection{1.2 Perché limiti e colimiti sono il test più severo}

Limiti e colimiti sono il cuore della ``universalità'' in category theory. Se OCT non riesce a trattarli in modo coerente, l'estensione resta locale. Se li tratta in modo formalmente stabile, l'estensione diventa sistemica.

Questo capitolo lavora esattamente su quel confine:

\begin{enumerate}
\def\labelenumi{\arabic{enumi}.}
\tightlist
\item
  mostra quando una costruzione universale classica resta universalmente valida in senso ordinativo;
\item
  mostra quando la costruzione resta formalmente lecita ma è ordinativamente vuota.
\end{enumerate}

\begin{center}\rule{0.5\linewidth}{0.5pt}\end{center}

\section{2. Tesi del capitolo}

La tesi è articolata in nove enunciati.

\begin{enumerate}
\def\labelenumi{\arabic{enumi}.}
\tightlist
\item
  L'isomorfismo classico resta invariato come nozione strutturale, ma non implica automaticamente isomorfismo ordinativo.
\item
  L'equivalenza classica di categorie non implica equivalenza ordinativa piena.
\item
  Limiti classici e limiti ordinativi coincidono solo sotto vincoli di coerenza ed emergenza.
\item
  Colimiti classici e colimiti ordinativi coincidono solo se l'aggregazione non annulla la funzione emergente.
\item
  È possibile definire nozioni di isomorfismo/equivalenza/limite/colimite ordinativi senza contraddire il quadro classico.
\item
  La non-ridondanza di OCT emerge in modo naturale proprio sui costrutti universali.
\item
  La selettività ordinativa è compatibile con il funtore di oblio verso la teoria classica.
\item
  La distinzione tra ``universale formale'' e ``universale reale'' è metodologicamente necessaria.
\item
  Questo capitolo prepara direttamente funtori, naturali e aggiunzioni del Capitolo 6.
\end{enumerate}

\begin{center}\rule{0.5\linewidth}{0.5pt}\end{center}

\section{3. Definizioni canoniche}

\subsection{Definizione 5.1 (Isomorfismo ordinativo)}

Siano \texttt{A\_Omega}, \texttt{B\_Omega} oggetti ordinativi. Un isomorfismo ordinativo è una coppia \texttt{f:\ A\_Omega-\textgreater{}B\_Omega}, \texttt{g:\ B\_Omega-\textgreater{}A\_Omega} tale che:

\begin{enumerate}
\def\labelenumi{\arabic{enumi}.}
\tightlist
\item
  \texttt{g\ o\ f\ =\ id\_A} e \texttt{f\ o\ g\ =\ id\_B} in senso classico;
\item
  \texttt{Adm\_Omega(f)}, \texttt{Adm\_Omega(g)};
\item
  le identità risultanti sono ordinativamente neutre.
\end{enumerate}

\subsection{Definizione 5.2 (Equivalenza ordinativa di categorie)}

Due categorie ordinative \texttt{C\_Omega}, \texttt{D\_Omega} sono ordinativamente equivalenti se esistono funtori \texttt{F}, \texttt{G} che realizzano equivalenza classica e preservano invarianti ordinativi dichiarati entro tolleranza.

\subsection{Definizione 5.3 (Limite ordinativo)}

Dato un diagramma \texttt{D}, un limite ordinativo è un limite classico \texttt{Lim(D)} che soddisfa:

\begin{enumerate}
\def\labelenumi{\arabic{enumi}.}
\tightlist
\item
  \texttt{Coh\_Omega(Lim(D))\ \textgreater{}=\ tau\_lim};
\item
  \texttt{Phi\_Omega(Lim(D))\ !=\ 0};
\item
  stabilità locale su vicinato contestuale dichiarato.
\end{enumerate}

\subsection{Definizione 5.4 (Colimite ordinativo)}

Dato \texttt{D}, un colimite ordinativo è un colimite classico \texttt{Colim(D)} che soddisfa:

\begin{enumerate}
\def\labelenumi{\arabic{enumi}.}
\tightlist
\item
  \texttt{Coh\_Omega(Colim(D))\ \textgreater{}=\ tau\_colim};
\item
  \texttt{Phi\_Omega(Colim(D))\ !=\ 0};
\item
  assenza di perdita emergenziale critica in aggregazione.
\end{enumerate}

\subsection{Definizione 5.5 (Universalità selettiva)}

Una costruzione universale è selettivamente valida in OCT se:

\begin{enumerate}
\def\labelenumi{\arabic{enumi}.}
\tightlist
\item
  è universale nel senso classico;
\item
  supera il gate ordinativo nel contesto osservativo.
\end{enumerate}

\subsection{Definizione 5.6 (Falso universale ordinativo)}

Costrutto formalmente universale ma ordinativamente invalido (sterile o degenerativo) nel dominio dichiarato.

\subsection{Definizione 5.7 (Scarto di equivalenza)}

Scarto tra equivalenza classica e ordinativa:

\texttt{Gap\_eq(Omega)\ :=\ Eq\_class\ -\ Eq\_ord,Omega}

inteso come divergenza di classificazione su un insieme di test strutturali.

\begin{center}\rule{0.5\linewidth}{0.5pt}\end{center}

\section{4. Sviluppo formale}

\subsection{4.1 Isomorfismo: conservazione e vincolo}

La nozione classica di isomorfismo resta intatta:

\begin{enumerate}
\def\labelenumi{\arabic{enumi}.}
\tightlist
\item
  esistenza di inversi;
\item
  composizione con identità.
\end{enumerate}

In OCT aggiungiamo:

\begin{enumerate}
\def\labelenumi{\arabic{enumi}.}
\tightlist
\item
  ammissibilità ordinativa dei morfismi coinvolti;
\item
  neutralità ordinativa delle identità risultanti.
\end{enumerate}

Ne segue una distinzione netta:

\begin{enumerate}
\def\labelenumi{\arabic{enumi}.}
\tightlist
\item
  isomorfismo classico puro;
\item
  isomorfismo classico + ordinativo.
\end{enumerate}

Errore evitato:

\begin{enumerate}
\def\labelenumi{\arabic{enumi}.}
\tightlist
\item
  inferire uguaglianza operativa totale da sola invertibilità strutturale.
\end{enumerate}

\subsection{4.2 Equivalenza di categorie: dove nasce lo scarto}

Nel classico, equivalenza categoriale significa ``stessa teoria fino a equivalenza''. In OCT questa affermazione resta vera sul piano strutturale, ma non è automaticamente vera sul piano ordinativo.

Perché può emergere scarto:

\begin{enumerate}
\def\labelenumi{\arabic{enumi}.}
\tightlist
\item
  \texttt{F} e \texttt{G} possono preservare forme ma non robustezza contestuale;
\item
  la stessa struttura può produrre dinamiche emergenti diverse in contesti differenti;
\item
  alcune trasformazioni possono introdurre perdita ontologica (D03).
\end{enumerate}

Quindi:

\begin{enumerate}
\def\labelenumi{\arabic{enumi}.}
\tightlist
\item
  \texttt{C\ \ \textasciitilde{}=\ \ D} classico non implica \texttt{C\ \ \textasciitilde{}=\ ord\ D};
\item
  serve controllo invarianti ordinativi sulle immagini funtoriali rilevanti.
\end{enumerate}

\subsection{4.3 Limiti: universalità non basta}

Nel classico, \texttt{Lim(D)} è definito dalla proprietà universale del cono terminale. In OCT questa proprietà resta necessaria ma non sufficiente.

Condizione aggiuntiva:

\begin{enumerate}
\def\labelenumi{\arabic{enumi}.}
\tightlist
\item
  il limite deve essere ordinativamente non sterile e non degenerativo.
\end{enumerate}

Caso tipico di falso universale:

\begin{enumerate}
\def\labelenumi{\arabic{enumi}.}
\tightlist
\item
  il limite classico esiste;
\item
  struttura formalmente corretta;
\item
  emergenza funzionale nulla (\texttt{Phi=0}) o margine coerentivo insufficiente.
\end{enumerate}

Conclusione:

\begin{enumerate}
\def\labelenumi{\arabic{enumi}.}
\tightlist
\item
  limite classico non implica limite ordinativo reale.
\end{enumerate}

\subsection{4.4 Colimiti: aggregare non significa integrare}

Il colimite classico privilegia la compatibilità di coconi e l'universalità di aggregazione. In OCT serve anche:

\begin{enumerate}
\def\labelenumi{\arabic{enumi}.}
\tightlist
\item
  conservazione di coerenza oltre soglia;
\item
  funzione emergente non nulla dell'aggregato.
\end{enumerate}

Differenza pratica:

\begin{enumerate}
\def\labelenumi{\arabic{enumi}.}
\tightlist
\item
  un colimite può ``unire'' formalmente componenti senza produrre integrazione ordinativa;
\item
  l'aggregazione può perfino aumentare rumore o annullare funzione.
\end{enumerate}

Questa è la ragione per cui parliamo di universalità selettiva:

\begin{enumerate}
\def\labelenumi{\arabic{enumi}.}
\tightlist
\item
  non ogni ``migliore aggregatore formale'' è ``migliore aggregatore reale''.
\end{enumerate}

\subsection{4.5 Relazione tra isomorfismi/equivalenze e limiti/colimiti}

Le quattro nozioni sono spesso studiate separatamente, ma in OCT sono interdipendenti.

\begin{enumerate}
\def\labelenumi{\arabic{enumi}.}
\tightlist
\item
  un forte \texttt{Gap\_eq} influenza trasferibilità di limiti/colimiti tra categorie equivalenti classiche;
\item
  limiti/colimiti ordinativamente instabili possono produrre falsi segnali di equivalenza operativa;
\item
  isomorfismi locali non garantiscono universalità ordinativa globale.
\end{enumerate}

Da qui segue un principio metodologico:

\begin{enumerate}
\def\labelenumi{\arabic{enumi}.}
\tightlist
\item
  equivalenze e universalità devono essere valutate in coppia quando il dominio è dinamico o ad alta sensibilità relazionale.
\end{enumerate}

\subsection{4.6 Criterio di accettazione per costrutti universali}

Per uso operativo, proponiamo una regola minima:

un limite/colimite può essere dichiarato ordinativamente valido solo se:

\begin{enumerate}
\def\labelenumi{\arabic{enumi}.}
\tightlist
\item
  proprietà universale classica verificata;
\item
  gate \texttt{Coh/Phi} superato;
\item
  robustezza minima su perturbazioni dichiarate;
\item
  classificazione non ricade in sterilità.
\end{enumerate}

Questa regola riduce sia overclaiming teorico sia regressioni applicative.

\subsection{4.7 Stabilità di equivalenza sotto variazione contestuale}

Per equivalenze classiche candidate a equivalenza ordinativa, va testata la stabilità su \texttt{N\_epsilon(Omega)}.

Richiesta minima:

\begin{enumerate}
\def\labelenumi{\arabic{enumi}.}
\tightlist
\item
  immagini funtoriali critiche non devono attraversare sistematicamente il confine A/B/C;
\item
  la divergenza residua deve restare entro banda dichiarata.
\end{enumerate}

In assenza di questo controllo:

\begin{enumerate}
\def\labelenumi{\arabic{enumi}.}
\tightlist
\item
  l'equivalenza ordinativa rischia di essere artefatto locale.
\end{enumerate}

\subsection{4.8 Matrice di fallimento specifica del capitolo}

Tipi di fallimento frequenti:

\begin{enumerate}
\def\labelenumi{\arabic{enumi}.}
\tightlist
\item
  \textbf{F-Iso1}: isomorfismo classico con perdita ordinativa su inverso;
\item
  \textbf{F-Eq1}: equivalenza classica con divergenza stabile di invarianti;
\item
  \textbf{F-Lim1}: limite classico sterile (\texttt{Phi=0});
\item
  \textbf{F-Lim2}: limite classico fragile (\texttt{Coh} al bordo);
\item
  \textbf{F-Col1}: colimite aggregativo con perdita emergenziale;
\item
  \textbf{F-Col2}: colimite formalmente corretto ma non robusto su perturbazioni.
\end{enumerate}

Questa matrice serve a collegare teoria e debugging nei capitoli successivi.

\subsection{4.9 Compatibilità con il recupero classico}

Il capitolo è coerente con F10:

\begin{enumerate}
\def\labelenumi{\arabic{enumi}.}
\tightlist
\item
  disattivando il filtro ordinativo, si recuperano isomorfismi/equivalenze/limiti/colimiti classici;
\item
  attivando il filtro, si ottiene selettività senza contraddire la struttura classica.
\end{enumerate}

Questo punto è essenziale per mantenere continuità didattica e matematica.

\subsection{4.10 Implicazione strategica per il libro}

La selettività introdotta qui modifica il modo di leggere l'intero corpus classico:

\begin{enumerate}
\def\labelenumi{\arabic{enumi}.}
\tightlist
\item
  non basta chiedere ``esiste?'';
\item
  bisogna chiedere ``esiste ed è ordinativamente reale nel contesto dichiarato?''.
\end{enumerate}

Questa doppia domanda diventa standard anche per il Capitolo 6 (funtori, naturali, aggiunzioni, monadi/comonadi).

\subsection{4.11 Condizione di non-ridondanza del capitolo}

Il capitolo è non ridondante se esiste almeno un caso per ciascuna coppia:

\begin{enumerate}
\def\labelenumi{\arabic{enumi}.}
\tightlist
\item
  isomorfismo classico valido / isomorfismo ordinativo invalido;
\item
  equivalenza classica valida / equivalenza ordinativa invalida;
\item
  limite classico esistente / limite ordinativo non reale;
\item
  colimite classico esistente / colimite ordinativo non reale.
\end{enumerate}

Questa condizione è soddisfatta dal programma teorematico F01, F05, F06, F07 già fissato nei capitoli precedenti.

\subsection{4.12 Trasporto di equivalenza e perdita ordinativa}

Quando una equivalenza classica \texttt{F:\ C-\textgreater{}D} viene usata come trasporto di risultati, in OCT bisogna valutare il costo ordinativo del trasporto.

Definiamo una misura sintetica:

\texttt{Loss\_tr\_Omega(F)\ :=\ w1*DeltaCoh\ +\ w2*DeltaPhi\ +\ w3*Instab},

dove:

\begin{enumerate}
\def\labelenumi{\arabic{enumi}.}
\tightlist
\item
  \texttt{DeltaCoh} misura la variazione media di coerenza su oggetti/morfismi test;
\item
  \texttt{DeltaPhi} misura variazione emergenziale;
\item
  \texttt{Instab} misura frequenza di attraversamenti A/B/C dopo trasporto.
\end{enumerate}

Interpretazione:

\begin{enumerate}
\def\labelenumi{\arabic{enumi}.}
\tightlist
\item
  \texttt{Loss\_tr} basso: trasporto ordinativamente affidabile;
\item
  \texttt{Loss\_tr} alto: equivalenza classica utile formalmente ma fragile operativamente.
\end{enumerate}

Questo blocco non sostituisce la teoria dell'equivalenza, ma aiuta a decidere quando usare o meno un trasporto come argomento forte in un dominio applicativo.

\subsection{4.13 Universalità sotto rumore strutturale}

Un limite o colimite può essere accettabile in condizioni ideali ma non in condizioni realistiche. Per questo proponiamo una verifica di universalità sotto rumore controllato.

Schema:

\begin{enumerate}
\def\labelenumi{\arabic{enumi}.}
\tightlist
\item
  calcolare il costrutto universale classico;
\item
  applicare perturbazioni compatibili con \texttt{N\_epsilon(Omega)};
\item
  verificare persistenza del gate ordinativo.
\end{enumerate}

Criterio:

\texttt{UnivRob\_Omega(U)\ =\ true} se per ogni perturbazione ammessa il costrutto resta nel dominio ordinativamente valido.

Conseguenze:

\begin{enumerate}
\def\labelenumi{\arabic{enumi}.}
\tightlist
\item
  distingue universalità teorica da universalità robusta;
\item
  riduce rischio di sovrastima su casi ``puliti'' ma non replicabili.
\end{enumerate}

\subsection{4.14 Distinzione tra limite informativo e limite operativo}

Nel trattamento di limiti e colimiti conviene separare:

\begin{enumerate}
\def\labelenumi{\arabic{enumi}.}
\tightlist
\item
  \textbf{limite informativo}: minimizza ambiguità descrittiva nel diagramma;
\item
  \textbf{limite operativo}: massimizza tenuta ordinativa nelle condizioni d'uso.
\end{enumerate}

In alcuni domini i due concetti coincidono, in altri divergono. La divergenza spiega perché un costrutto formalmente ottimale possa risultare pragmaticamente debole.

Questo punto è importante per evitare un errore frequente:

\begin{enumerate}
\def\labelenumi{\arabic{enumi}.}
\tightlist
\item
  usare criteri informativi puri per giustificare decisioni operative ad alta criticità.
\end{enumerate}

\subsection{4.15 Pattern tipici di falsa equivalenza}

Osserviamo quattro pattern ricorrenti.

\begin{enumerate}
\def\labelenumi{\arabic{enumi}.}
\item
  \textbf{Pattern E1 - Equivalenza di forma, divergenza di traiettoria}\\
  stesse relazioni iniziali, dinamiche emergenti divergenti.
\item
  \textbf{Pattern E2 - Equivalenza locale, collasso globale}\\
  traduzioni corrette modulo per modulo, ma incompatibilità cumulativa.
\item
  \textbf{Pattern E3 - Universalità sterile}\\
  soluzione universale formale senza funzione emergente apprezzabile.
\item
  \textbf{Pattern E4 - Aggregazione degenerativa}\\
  colimite che unifica ma riduce capacità operativa del sistema aggregato.
\end{enumerate}

Riconoscere questi pattern accelera sia proof design sia debugging sperimentale.

\subsection{4.16 Protocollo minimo di verifica per limiti e colimiti}

Per rendere replicabile la valutazione ordinativa proponiamo una pipeline minima in sei passi.

\begin{enumerate}
\def\labelenumi{\arabic{enumi}.}
\tightlist
\item
  costruzione del costrutto classico (\texttt{Lim} o \texttt{Colim}) con prova formale standard;
\item
  specifica del contesto \texttt{Omega} e delle soglie \texttt{tau} pertinenti;
\item
  misura puntuale di \texttt{Coh\_Omega} e \texttt{Phi\_Omega};
\item
  test su vicinato \texttt{N\_epsilon(Omega)} con perturbazioni dichiarate;
\item
  classificazione finale (\texttt{accepted}, \texttt{revise}, \texttt{reject});
\item
  registrazione nel claim ledger con evidenza \texttt{S0-S3}.
\end{enumerate}

Questa pipeline crea continuità tra matematica e metodo:

\begin{enumerate}
\def\labelenumi{\arabic{enumi}.}
\tightlist
\item
  la parte 1 garantisce rigore classico;
\item
  le parti 2-4 garantiscono rigore ordinativo;
\item
  le parti 5-6 garantiscono rigore epistemico.
\end{enumerate}

In assenza di uno di questi passaggi, la decisione resta incompleta: o formalmente robusta ma operativamente cieca, o operativamente intuitiva ma teoricamente debole.

\begin{center}\rule{0.5\linewidth}{0.5pt}\end{center}

\section{5. Proposizioni e teoremi locali}

\subsection{Proposizione 5.1 (Necessità del vincolo ordinativo sugli isomorfismi)}

Enunciato:
L'invertibilità classica non è condizione sufficiente per isomorfismo ordinativo.

Intuizione:
gli inversi possono essere tipologicamente corretti ma emergenzialmente sterili o non robusti.

Stato: \texttt{validated}.

\subsection{Proposizione 5.2 (Universalità selettiva dei limiti)}

Enunciato:
Esistono limiti classici non accettabili come limiti ordinativi nel medesimo \texttt{Omega}.

Intuizione:
proprietà universale formale e validità ordinativa sono livelli distinti.

Stato: \texttt{in\_proof}.

\subsection{Proposizione 5.3 (Universalità selettiva dei colimiti)}

Enunciato:
Esistono colimiti classici con perdita emergenziale critica, quindi non ordinativamente validi.

Intuizione:
aggregazione strutturale non implica integrazione funzionale.

Stato: \texttt{in\_proof}.

\subsection{Teorema 5.4 (Scarto di equivalenza non nullo in domini dinamici)}

Ipotesi:

\begin{enumerate}
\def\labelenumi{\arabic{enumi}.}
\tightlist
\item
  coppia di categorie classiche equivalenti;
\item
  mapping ordinativo su oggetti/morfismi critici;
\item
  divergenza stabile su almeno un invariante ordinativo.
\end{enumerate}

Tesi:
\texttt{Gap\_eq(Omega)\ \textgreater{}\ 0}.

Proof sketch:

\begin{enumerate}
\def\labelenumi{\arabic{enumi}.}
\tightlist
\item
  l'equivalenza classica garantisce traduzione strutturale;
\item
  gli invarianti ordinativi rilevano differenze funzionali/robuste non assorbite dalla traduzione;
\item
  la divergenza stabile implica scarto positivo.
\end{enumerate}

Conseguenze:
dimostra formalmente il valore aggiunto di OCT a livello categoriale globale.

Stato: \texttt{in\_proof}.

\subsection{Teorema 5.5 (Criterio minimo per universalità ordinativa)}

Ipotesi:

\begin{enumerate}
\def\labelenumi{\arabic{enumi}.}
\tightlist
\item
  costrutto universale classico \texttt{U(D)} (limite o colimite);
\item
  gate ordinativo soddisfatto su \texttt{U(D)};
\item
  robustezza locale su \texttt{N\_epsilon(Omega)}.
\end{enumerate}

Tesi:
\texttt{U(D)} è universalmente accettabile in senso ordinativo nel dominio dichiarato.

Proof sketch:

\begin{enumerate}
\def\labelenumi{\arabic{enumi}.}
\tightlist
\item
  la proprietà universale classica stabilisce correttezza strutturale;
\item
  gate e robustezza stabiliscono realtà compositiva;
\item
  l'unione dei due livelli produce universalità ordinativa.
\end{enumerate}

Conseguenze:
fornisce una regola unificata per trattare limiti e colimiti in OCT.

Stato: \texttt{in\_proof}.

\subsection{Proposizione 5.6 (Non-invarianza automatica del trasporto funtoriale)}

Enunciato:
Un funtore che realizza equivalenza classica non preserva automaticamente il livello di robustezza ordinativa.

Intuizione:
la preservazione strutturale non implica preservazione del margine \texttt{m\_Omega}.

Stato: \texttt{in\_review}.

\subsection{Teorema 5.7 (Vincolo di robustezza universale)}

Ipotesi:

\begin{enumerate}
\def\labelenumi{\arabic{enumi}.}
\tightlist
\item
  esistenza classica di \texttt{U(D)} (limite o colimite);
\item
  validità ordinativa puntuale di \texttt{U(D)} in \texttt{Omega};
\item
  \texttt{UnivRob\_Omega(U)=true}.
\end{enumerate}

Tesi:
\texttt{U(D)} è universalmente utilizzabile in modo ordinativamente affidabile nel dominio dichiarato.

Proof sketch:

\begin{enumerate}
\def\labelenumi{\arabic{enumi}.}
\tightlist
\item
  la validità puntuale garantisce ammissibilità iniziale;
\item
  la robustezza su vicinato evita falsi positivi da fitting locale;
\item
  quindi il costrutto è idoneo a uso operativo ripetibile.
\end{enumerate}

Conseguenze:
fornisce criterio più forte del solo gate puntuale per capitoli applicativi futuri.

Stato: \texttt{in\_proof}.

\begin{center}\rule{0.5\linewidth}{0.5pt}\end{center}

\section{6. Implicazioni operative}

\subsection{6.1 Per progettazione matematica}

Ogni volta che si usa ``isomorfo'', ``equivalente'', ``limite'', ``colimite'' nel manoscritto OCT, è raccomandato specificare:

\begin{enumerate}
\def\labelenumi{\arabic{enumi}.}
\tightlist
\item
  livello classico;
\item
  livello ordinativo;
\item
  eventuale gap tra i due livelli.
\end{enumerate}

\subsection{6.2 Per benchmarking}

I benchmark dovrebbero includere set specifici:

\begin{enumerate}
\def\labelenumi{\arabic{enumi}.}
\tightlist
\item
  casi con equivalenza classica ma divergenza ordinativa;
\item
  casi con universalità classica ma sterilità emergenziale;
\item
  casi robusti in cui classico e ordinativo coincidono.
\end{enumerate}

Questo evita valutazioni sbilanciate su soli casi favorevoli.

\subsection{6.3 Per interpretazione dei risultati}

Quando un costrutto fallisce ordinativamente, non significa che la teoria classica sia ``sbagliata''. Significa che, nel dominio osservato, il livello classico non basta per decisioni operative.

\subsection{6.4 Per pipeline modulari}

Nei sistemi composti da più moduli:

\begin{enumerate}
\def\labelenumi{\arabic{enumi}.}
\tightlist
\item
  limiti e colimiti vanno valutati non solo per correttezza formale;
\item
  vanno valutati anche per effetti su coerenza globale e funzione emergente.
\end{enumerate}

\subsection{6.5 Per il Capitolo 6}

Funtori e naturali del Capitolo 6 erediteranno direttamente questi vincoli:

\begin{enumerate}
\def\labelenumi{\arabic{enumi}.}
\tightlist
\item
  preservazione classica;
\item
  preservazione ordinativa;
\item
  tracciamento esplicito degli scarti.
\end{enumerate}

\subsection{6.6 Per audit di equivalenze e universalità}

Nelle revisioni tecniche è utile allegare una tabella per ogni costrutto candidato:

\begin{enumerate}
\def\labelenumi{\arabic{enumi}.}
\tightlist
\item
  tipo (\texttt{iso}, \texttt{eq}, \texttt{lim}, \texttt{colim});
\item
  validità classica (\texttt{si/no});
\item
  validità ordinativa puntuale (\texttt{si/no});
\item
  robustezza su \texttt{N\_epsilon(Omega)} (\texttt{si/no});
\item
  stato finale (\texttt{accepted}, \texttt{revise}, \texttt{reject}).
\end{enumerate}

Questa tabella rende comparabili team diversi e riduce ambiguità nei passaggi tra release.

\subsection{6.7 Per continuità editoriale del manoscritto}

Nel testo finale dell'opera è consigliato mantenere per ciascun costrutto universale un box standard:

\begin{enumerate}
\def\labelenumi{\arabic{enumi}.}
\tightlist
\item
  definizione classica sintetica;
\item
  estensione ordinativa;
\item
  controesempio tipico;
\item
  criterio di accettazione.
\end{enumerate}

In questo modo il lettore non specialista mantiene orientamento, mentre il lettore tecnico trova rapidamente i punti di controllo formale.

\subsection{6.8 Priorità operative per i prossimi capitoli}

Per massimizzare continuità con il Capitolo 6 è utile fissare tre priorità immediate:

\begin{enumerate}
\def\labelenumi{\arabic{enumi}.}
\tightlist
\item
  costruire esempi in cui un funtore preserva struttura ma non preserva classe A/B/C;
\item
  esplicitare almeno una aggiunzione classica che richiede vincolo ordinativo aggiuntivo;
\item
  preparare un caso monadico in cui la composizione iterata mostra differenza tra validità locale e validità globale.
\end{enumerate}

Queste priorità evitano discontinuità tra capitoli e permettono di riusare direttamente la grammatica del Capitolo 5 nel trattamento di naturali, aggiunzioni e monadi/comonadi.

\subsection{6.9 Regola pratica di comunicazione scientifica}

Quando si presenta un risultato su limiti o colimiti in contesto OCT, è utile dichiarare sempre in apertura:

\begin{enumerate}
\def\labelenumi{\arabic{enumi}.}
\tightlist
\item
  cosa è dimostrato in senso classico;
\item
  cosa è accettato in senso ordinativo;
\item
  quale parte resta in stato \texttt{in\_proof} o \texttt{revise}.
\end{enumerate}

Questa regola riduce incomprensioni tra lettori con background diversi e migliora la qualità delle revisioni tra pari, perché separa subito livello matematico, livello operativo e livello di evidenza disponibile.
Inoltre facilità la produzione di abstract tecnici sintetici coerenti con il claim ledger, utili per repository, preprint e comunicazioni istituzionali.

\begin{center}\rule{0.5\linewidth}{0.5pt}\end{center}

\section{7. Limiti e non-claim}

Limiti espliciti:

\begin{enumerate}
\def\labelenumi{\arabic{enumi}.}
\tightlist
\item
  il capitolo non chiude tutte le prove complete di F05/F06/F07;
\item
  la robustezza di limiti/colimiti resta dipendente da modelli perturbativi dichiarati;
\item
  la misura di \texttt{Gap\_eq} è presentata in forma concettuale minima, da raffinare.
\end{enumerate}

Failure mode riconosciuti:

\begin{enumerate}
\def\labelenumi{\arabic{enumi}.}
\tightlist
\item
  considerare \texttt{Gap\_eq} nullo perché non misurato correttamente;
\item
  trattare casi sterilità come eccezioni trascurabili senza analisi;
\item
  usare solo esempi in cui classico e ordinativo coincidono;
\item
  applicare equivalenze categoriali senza verifica su invarianti ordinativi critici.
\end{enumerate}

Box C - Non-claim

Questo capitolo non implica:

\begin{enumerate}
\def\labelenumi{\arabic{enumi}.}
\tightlist
\item
  che i costrutti classici perdano validità matematica intrinseca;
\item
  che ogni limite o colimite classico sia ordinativamente invalido;
\item
  che OCT elimini il bisogno di dimostrazione classica dei costrutti universali.
\end{enumerate}

\begin{center}\rule{0.5\linewidth}{0.5pt}\end{center}

\section{8. Claim Ledger (S0-S3)}

{\def\LTcaptype{none} % do not increment counter
\begin{longtable}[]{@{}lllll@{}}
\toprule\noalign{}
ID & Claim & Tipo & Evidenza & Stato \\
\midrule\noalign{}
\endhead
\bottomrule\noalign{}
\endlastfoot
C5.1 & Isomorfismo classico e isomorfismo ordinativo non coincidono in generale & formale & S2 & in\_proof \\
C5.2 & Equivalenza classica non implica equivalenza ordinativa piena & formale & S2 & in\_proof \\
C5.3 & Limiti classici richiedono filtro ordinativo per essere ``reali'' in \texttt{Omega} & formale & S2 & in\_proof \\
C5.4 & Colimiti classici possono essere aggregativi ma non integrativi & formale-operativo & S2 & in\_proof \\
C5.5 & La nozione di universalità selettiva è compatibile con il recupero classico & metateorico & S1 & validated \\
C5.6 & La matrice F-Iso/F-Eq/F-Lim/F-Col migliora la diagnosi operativa & metodologico & S1 & in\_review \\
\end{longtable}
}

\begin{center}\rule{0.5\linewidth}{0.5pt}\end{center}

\section{9. Bridge al Capitolo 6}

Il Capitolo 5 ha esteso in senso ordinativo i costrutti universali principali del quadro classico:

\begin{enumerate}
\def\labelenumi{\arabic{enumi}.}
\tightlist
\item
  isomorfismi;
\item
  equivalenze;
\item
  limiti;
\item
  colimiti.
\end{enumerate}

Il Capitolo 6 affronterà ora il livello trasformazionale superiore:

\begin{enumerate}
\def\labelenumi{\arabic{enumi}.}
\tightlist
\item
  funtori e trasformazioni naturali;
\item
  aggiunzioni;
\item
  monadi e comonadi.
\end{enumerate}

In particolare, il nodo centrale sarà capire quando una trasformazione tra categorie preserva non solo la forma, ma anche la validità ordinativa delle strutture trasportate.

\begin{center}\rule{0.5\linewidth}{0.5pt}\end{center}

\section{Riferimenti}

Riferimenti di framework interno:

\begin{enumerate}
\def\labelenumi{\arabic{enumi}.}
\tightlist
\item
  Ghioni, F. (2026). \texttt{TE\_CORE\_v5.1.md}.
\item
  Ghioni, F. (2026). \texttt{Ordinative\_Set\_Theory\_OST\_A\_Concise\_Guide\_For\_AI\_v2\_1.md}.
\item
  \texttt{OCT\_FOUNDATIONAL\_BOOK\_CHAPTER\_04\_v1\_0.md}.
\item
  \texttt{OCT\_FOUNDATIONAL\_BOOK\_CHAPTER\_08\_v1\_0.md}.
\item
  \texttt{OCT\_FOUNDATIONAL\_BOOK\_CHAPTER\_09\_v1\_0.md}.
\item
  \texttt{OCT\_FOUNDATIONAL\_BOOK\_CHAPTER\_10\_v1\_0.md}.
\end{enumerate}

Riferimenti primari:

\begin{enumerate}
\def\labelenumi{\arabic{enumi}.}
\tightlist
\item
  Mac Lane, S. (1998). \emph{Categories for the Working Mathematician}.
\item
  Riehl, E. (2016). \emph{Category Theory in Context}.
\item
  Spivak, D. (2014). \emph{Category Theory for the Sciences}.
\end{enumerate}
